\section{An exact attainment criterion and the twelve-option certificate}
\label{app:attainment}
For $n\ge4$, let $D=n^2-3n+1$ and let $\mathcal C$ be the set of
derangement cycle types in $S_n$. For a Latin square $T$, draw two
distinct rows $x,y$ uniformly and put $q_T(C)=\Pr(x^{-1}y\in C)$.
Let $Q_n$ be the convex hull of these histograms over Latin squares of
order $n$. Repeated histograms need only be included once.

\paragraph{Cycle-type criterion.}
The bound $L_n$ is attained simultaneously for real scores and
correctness differences if and only if there exists $q\in Q_n$ with
\begin{equation}
 \sum_{C\in\mathcal C} q(C)\chi_\lambda(C)
 \ \le\ f_\lambda/D \qquad(\lambda\vdash n,\ \lambda\ne(n)).
 \label{eq:attain-lp}
\end{equation}
These are rational linear feasibility constraints: write
$q=\sum_T w_Tq_T$, $w_T\ge0$, $\sum_Tw_T=1$.
Whenever feasible, at most $|\mathcal C|$ templates with rational
weights suffice. In particular this test uses cycle types of row
pairs rather than all $n!$ coordinates of a response. A restricted
list suffices to prove attainment; an infeasibility claim requires
control of all eligible histograms.

\emph{Proof.} Average a feasible design independently over option
and slot relabelings. This preserves the uniform marginals. Each
relabeling merely permutes the coordinates of a profile, so the
worst-case risk cannot increase in either profile class. After
averaging, the design selects a template and applies two uniform
relabelings. Conversely, every point in $Q_n$ is realized in this way.
The nonconstant eigenvalues in Eq.~\eqref{eq:character-certificate}
are $1/n+(n-1)\sum_C q(C)\chi_\lambda(C)/(nf_\lambda)$.
Comparing these values with $L_n=1/n+(n-1)/(nD)$ gives
Eq.~\eqref{eq:attain-lp} in both directions for real profiles.
The lower bound already has a correctness witness, so a feasible
solution is sharp for both classes. A rational feasible polytope
has a rational vertex. Its histogram lies in a simplex spanned by
affinely independent template histograms, at most $|\mathcal C|$
of them since their entries sum to one. Its barycentric coordinates
in that simplex are rational.

The dual form gives a certificate of impossibility. Define
$b_{T\lambda}=\sum_Cq_T(C)\chi_\lambda(C)/f_\lambda-1/D$ for
nontrivial $\lambda$. The criterion fails precisely when there is a
probability vector $u$ on these partitions such that
$\sum_\lambda u_\lambda b_{T\lambda}>0$ for every template.
Indeed the convex hull of the finite row vectors $b_T$ then misses
the closed negative orthant. Strict separation supplies a nonzero
nonnegative normal (otherwise its supremum on the orthant is infinite);
normalizing its sum gives $u$. Conversely such a $u$ excludes every
feasible mixture.

\paragraph{What balance is necessary?}
The two dyad inequalities have respective left sides $t$ and $1-t$,
where $t=\sum_Cq(C)m_2(C)$ and $m_2(C)$ counts transpositions in the
cycle decomposition. Their right sides are $s/D$ and $a/D$ with
$s+a=D$. Hence both must be equalities. This gives an initial scalar
test before the other inequalities. The same equality is necessary
even for correctness alone by Eq.~\eqref{eq:swaps}. For an exact
criterion for that class alone, replace $\chi_\lambda(C)/f_\lambda$
in the normalized inequalities by
\[
 h_z(C)=\mathbb E_{x,g}[z(x)z(xg)]/\mathbb E_x[z(x)^2]
\]
for each centered nonconstant ternary profile $z$ (center after
choosing the ternary entries), with independent uniform $x\in S_n$
and $g\in C$. Indeed the off-diagonal normalized second moment is
then $\sum_Cq(C)h_z(C)$; the diagonal contribution is unchanged.
These are again finitely many rational inequalities, though generally
far more numerous.

\paragraph{Screening the remaining degrees.}
For a candidate $q$, put $H(q)^2=\sum_C q(C)^2 n!/|C|$.
Column orthogonality gives
$\sum_\lambda(\sum_Cq(C)\chi_\lambda(C))^2=H(q)^2$.
Thus all components of degree at least $DH(q)$ satisfy the test
automatically. This bound can shorten an exact check without
assuming that the two dyad components are the worst ones.

\paragraph{Twelve options.}
The two squares in \texttt{certificates/position-templates.json}
under key \texttt{12} have total transposition counts 54 and 27 over
their 66 unordered row pairs. Draw the first square with probability
$23/109$ and the second with probability $86/109$, then shuffle
symbols and slots independently. The mean transposition count is
$(23\cdot54+86\cdot27)/(109\cdot66)=54/109=s/D$.
For this histogram $186<109H(q)<187$, so degrees at least 187 are
automatically covered. The accompanying exact certificate checks all
76 nontrivial partitions. The largest normalized risk is $10/109$
on $(10,2)$ and $(10,1,1)$; the largest remaining value is
$1873/21582<10/109$. One verifier uses the Murnaghan--Nakayama rule;
a second expands the Frobenius character formula and recomputes the
pair types from the rows. By the criterion and the correctness lower
bound, $10/109$ is the exact optimum for both classes.
